\documentclass[12pt,a4paper,oneside]{memoir}
\usepackage[utf8]{inputenc}
\usepackage[T1]{fontenc}
\usepackage{lmodern}
\usepackage{xurl}
\usepackage[spanish, es-tabla]{babel}
\addto\captionsspanish{
  \renewcommand{\listtablename}{Índice de Tablas}
  \renewcommand{\listfigurename}{Índice de Figuras}
}
\usepackage{amsmath,amssymb}
\usepackage{graphicx}
\usepackage{threeparttable}
\usepackage{pdflscape}
\usepackage{listings}
\lstset{
    language=Python,
    basicstyle=\ttfamily\footnotesize,
    keywordstyle=\color[HTML]{0000FF},
    commentstyle=\color[HTML]{008000},
    stringstyle=\color[HTML]{BA2121},
    numberstyle=\tiny\color[HTML]{888888},
    emph={True,False,None,self},
    emphstyle=\color[HTML]{008080},
    numbers=left,
    frame=single,
    rulecolor=\color[HTML]{CCCCCC},
    backgroundcolor=\color[HTML]{F9F9F9},
    lineskip=-1pt,
    aboveskip=5pt,
    belowskip=5pt,
    showstringspaces=false,
    breaklines=true,
    breakatwhitespace=false,
    postbreak=\mbox{\textcolor{gray}{$\hookrightarrow$}\space},
    inputencoding=utf8,
    extendedchars=true,
    literate=
        {á}{{\'a}}1 {é}{{\'e}}1 {í}{{\'i}}1
        {ó}{{\'o}}1 {ú}{{\'u}}1 {ñ}{{\~n}}1
        {Á}{{\'A}}1 {É}{{\'E}}1 {Í}{{\'I}}1
        {Ó}{{\'O}}1 {Ú}{{\'U}}1 {Ñ}{{\~N}}1
}
\nouppercaseheads
\renewcommand{\chaptermark}[1]{\markboth{\footnotesize\chaptername\ \thechapter.\ #1}{\footnotesize\chaptername\ \thechapter.\ #1}}
\usepackage{booktabs}
\usepackage{ragged2e}
\usepackage[
    protrusion=true,
    expansion=false
]{microtype}
\usepackage{booktabs}
\usepackage{physics}
\usepackage[version=4]{mhchem}

% Elimina numeración de secciones
\setcounter{secnumdepth}{1} %acá se utiliza para quitar la numeración desde: -1=chapter, 0=section, 1=subsection, etc.

% Pero mantiene todos los niveles en el índice
\setcounter{tocdepth}{3}
%sirven para enumerar en un unico continuo las imágenes, tablas y ecuaciones: (1,2,3,4,5,..., sucesivamente)
%\usepackage{chngcntr}
%\counterwithout{figure}{chapter}
%\counterwithout{table}{chapter}
%\counterwithout{equation}{chapter}

\usepackage{nameref}
\usepackage{gensymb}
\usepackage{atbegshi}
\usepackage{eso-pic}
\usepackage{float}
\usepackage{caption}
\usepackage{csquotes} 
\usepackage[absolute,overlay]{textpos}
\usepackage{array}
\usepackage[top=2.5cm, bottom=2.5cm, left=2.5cm, right=2.5cm]{geometry}
\usepackage{multirow}
\usepackage{tabulary}
\usepackage{subcaption}
\usepackage{tikz}
\usepackage[backend=biber,style=apa,sorting=ynt]{biblatex}

\defbibheading{bibliography}[\bibname]{%
    \chapter*{#1}%
    \markboth{#1}{#1}%
}

\renewcommand{\bibfont}{\footnotesize}

\AtEveryBibitem{\clearfield{title}}
\DeclareLanguageMapping{spanish}{spanish-apa}
\addbibresource{bibliografia.bib}
\usepackage[dvipsnames,svgnames]{xcolor}
\usepackage[colorinlistoftodos]{todonotes}
\usepackage{pgfgantt}
\definecolor{barblue}{RGB}{255,111,36}
\definecolor{groupcolor}{RGB}{255,184,148}
\definecolor{linkred}{RGB}{165,0,33}
\usepackage{imakeidx}
\usepackage{hyperref}
\hypersetup{
colorlinks,
citecolor=MidnightBlue,    % Muy oscuro
linkcolor=DarkBlue,        % Oscuro
urlcolor=DarkBlue,        % Medio
filecolor=DarkSlateBlue    % Oscuro
}
\makeindex[
    title=Índice Alfabético,
    columns=2,                          % 2 columnas
    columnsep=1cm,                  % Espacio entre columnas
    intoc]                              % Que aparezca en el índice general


\usepackage{titlesec}
\usepackage{acro}

% Primera mención: Español (en inglés English, SIGLA)
\NewAcroTemplate{long-short-es}{%
  \acroiffirstTF
    {\acrowrite{long}\acspace(\acroifT{foreign}{en~inglés~\acrowrite{foreign},~}\acrowrite{short})}
    {\acrowrite{short}}}
\acsetup{first-style = long-short-es}
\acsetup{format/foreign = \itshape}
\acsetup{make-links = true}

\definecolor{chaptercolor}{RGB}{0,0,0}
\titleformat{\chapter}[display]
{\normalfont\bfseries\color{chaptercolor}}
{\raggedright\Huge\MakeUppercase{\chaptertitlename}\ {\fontsize{40}{45}\selectfont\thechapter}} % Capítulo X
{20pt}
{\titlerule\vspace{10pt}\Huge\raggedleft} % Línea + Título
[\vspace{10pt}\titlerule] % Línea final
\titlespacing*{\chapter}{0pt}{50pt}{30pt}

\addto\captionsspanish{\renewcommand{\tablename}{Tabla}}

\captionnamefont{\normalfont\bfseries}    % "Tabla" en negrita
\captionstyle{\centering}                 % Centrado

\linespread{1.5}
\emergencystretch=2em


\title{Estudio Hidrodinámico de la Interacción entre Eyecciones de Masa Coronal Homólogas}
\author {Cristian David Chavez Aponte}
\date{\today}
\newcommand{\tutor}{Dr. Miguel Andrés Páez Murcia}
\newcommand{\universidad}{Universidad Pedagógica y Tecnológica de Colombia}
\newcommand{\facultad}{Facultad de Ciencias Básicas, Escuela de Física}
\newcommand{\programa}{Escuela de Física}
\newcommand{\lugar}{Tunja, Boyacá, Colombia}


\setlrmarginsandblock{3.5cm}{2.5cm}{*}
\setulmarginsandblock{2.5cm}{2.5cm}{*}
\checkandfixthelayout

\usepackage{transparent} % Permite control de opacidad con \transparent

% Marca de agua con opacidad
\newcommand\BackgroundPic{
  \put(0,0){
    \parbox[b][\paperheight]{\paperwidth}{
      \vfill
      \centering
      \transparent{0.05}
      \includegraphics[width=0.7\paperwidth]{Imágenes/escudo.png}
      \vfill
    }
  }
}
\newcommand{\margenIzquierdo}{
  \AddToShipoutPictureBG*{
    \ifnum\value{page}>2
      \begin{tikzpicture}[remember picture,overlay]
        \draw[white, thick, line width=2pt] (current page.south west) ++(2.55cm,5) -- ++ (0,25);
      \end{tikzpicture}
    \fi
  }
} 
\AtBeginShipout{\margenIzquierdo}


%================ACRÓNIMOS==================================

\DeclareAcronym{AIA}{short = AIA, long = {Conjunto de Imágenes Atmosféricas}, foreign = {Atmospheric Imaging Assembly}}
\DeclareAcronym{AR}{short = AR, long = {Región Activa}, foreign = {Active Region},
  short-plural-form = {ARs}, long-plural-form = {Regiones Activas}, foreign-plural-form = {Active Regions}}
\DeclareAcronym{AU}{short = AU, long = {Unidad Astronómica}, foreign = {Astronomical Unit}}
\DeclareAcronym{CACTus}{short = CACTus, long = {Rastreo de CMEs Asistido por Computadora}, foreign = {Computer Aided CME Tracking}}
\DeclareAcronym{CDAW}{short = CDAW, long = {Taller de Análisis de Datos Coordinados}, foreign = {Coordinated Data Analysis Workshop}}
\DeclareAcronym{CDS}{short = CDS, long = {Espectrómetro de Diagnóstico Coronal}, foreign = {Coronal Diagnostic Spectrometer}}
\DeclareAcronym{CH}{short = CH, long = {Agujero Coronal}, foreign = {Coronal Hole},
  short-plural-form = {CHs}, long-plural-form = {Agujeros Coronales}, foreign-plural-form = {Coronal Holes}}
\DeclareAcronym{CME}{short = CME, long = {Eyección de Masa Coronal}, foreign = {Coronal Mass Ejection},
  short-plural-form = {CMEs}, long-plural-form = {Eyecciones de Masa Coronal}, foreign-plural-form = {Coronal Mass Ejections}}
\DeclareAcronym{DKIST}{short = DKIST, long = {Telescopio Solar Daniel K. Inouye}, foreign = {Daniel K. Inouye Solar Telescope}}
\DeclareAcronym{DST}{short = DST, long = {Tiempo de Perturbación de Tormenta}, foreign = {Disturbance Storm Time}}
\DeclareAcronym{EFR}{short = EFR, long = {Región de Flujo Emergente}, foreign = {Emerging Flux Region},
  short-plural-form = {EFRs}, long-plural-form = {Regiones de Flujo Emergente}, foreign-plural-form = {Emerging Flux Regions}}
\DeclareAcronym{EIS}{short = EIS, long = {Espectrómetro de Imágenes en EUV}, foreign = {EUV Imaging Spectrometer}}
\DeclareAcronym{EUV}{short = EUV, long = {Ultravioleta Extremo}, foreign = {Extreme Ultraviolet}}
\DeclareAcronym{EUVI}{short = EUVI, long = {Generador de Imágenes en UV Extremo}, foreign = {EUV Imager}}
\DeclareAcronym{EVE}{short = EVE, long = {Experimento de Variabilidad en EUV}, foreign = {EUV Variability Experiment}}
\DeclareAcronym{FIELDS}{short = FIELDS, long = {Investigación de Campos Electromagnéticos}, foreign = {Electromagnetic Fields Investigation}}
\DeclareAcronym{HELCATS}{short = HELCATS, long = {Servicio de Catalogación, Análisis y Técnicas Heliosféricas}, foreign = {Heliospheric Cataloguing, Analysis and Techniques Service}}
\DeclareAcronym{HI}{short = HI, long = {Generador de Imágenes Heliosféricas}, foreign = {Heliospheric Imager}}
\DeclareAcronym{HMI}{short = HMI, long = {Generador de Imágenes Heliosismológicas y Magnéticas}, foreign = {Helioseismic and Magnetic Imager}}
\DeclareAcronym{ICME}{short = ICME, long = {CME Interplanetaria}, foreign = {Interplanetary CME},
  short-plural-form = {ICMEs}, long-plural-form = {Eyecciones de Masa Coronal Interplanetarias}, foreign-plural-form = {Interplanetary Coronal Mass Ejections}}
\DeclareAcronym{IMaX}{short = IMaX, long = {Experimento de Magnetógrafo de Imágenes}, foreign = {Imaging Magnetograph eXperiment}}
\DeclareAcronym{IMF}{short = IMF, long = {Campo Magnético Interplanetario}, foreign = {Interplanetary Magnetic Field}}
\DeclareAcronym{ISOIS}{short = ISOIS, long = {Investigación Científica Integrada del Sol}, foreign = {Integrated Science Investigation of the Sun}}
\DeclareAcronym{LASCO}{short = LASCO, long = {Coronógrafo Espectroscópico de Gran Ángulo}, foreign = {Large Angle and Spectroscopic Coronagraph}}
\DeclareAcronym{MC}{short = MC, long = {Nube Magnética}, foreign = {Magnetic Cloud},
  short-plural-form = {MCs}, long-plural-form = {Nubes Magnéticas}, foreign-plural-form = {Magnetic Clouds}}
\DeclareAcronym{MDI}{short = MDI, long = {Generador de Imágenes Doppler Michelson}, foreign = {Michelson Doppler Imager}}
\DeclareAcronym{MHD}{short = MHD, long = {Magnetohidrodinámica}, foreign = {Magnetohydrodynamics}}
\DeclareAcronym{NASA}{short = NASA, long = {Administración Nacional de Aeronáutica y del Espacio}, foreign = {National Aeronautics and Space Administration}}
\DeclareAcronym{NFI}{short = NFI, long = {Generador de Imágenes de Campo Estrecho}, foreign = {Narrow Field Imager}}
\DeclareAcronym{PFU}{short = PFU, long = {Unidades de Flujo de Partículas}, foreign = {Particle Flux Unit}}
\DeclareAcronym{PSP}{short = PSP, long = {Sonda Solar Parker}, foreign = {Parker Solar Probe}}
\DeclareAcronym{PUNCH}{short = PUNCH, long = {Polarímetro para Unificar la Corona y la Heliosfera}, foreign = {Polarimeter to UNify the Corona and Heliosphere}}
\DeclareAcronym{SDO}{short = SDO, long = {Observatorio de Dinámica Solar}, foreign = {Solar Dynamics Observatory}}
\DeclareAcronym{SECCHI}{short = SECCHI, long = {Investigación Coronal y Heliosférica de la Conexión Sol-Tierra}, foreign = {Sun Earth Connection Coronal and Heliospheric Investigation}}
\DeclareAcronym{SEP}{short = SEP, long = {Partícula Energética Solar}, foreign = {Solar Energetic Particle},
  short-plural-form = {SEPs}, long-plural-form = {Partículas Energéticas Solares}, foreign-plural-form = {Solar Energetic Particles}}
\DeclareAcronym{SF}{short = SF, long = {Llamarada Solar}, foreign = {Solar Flare},
  short-plural-form = {SFs}, long-plural-form = {Llamaradas Solares}, foreign-plural-form = {Solar Flares}}
\DeclareAcronym{SILSO}{short = SILSO, long = {Índice de Manchas Solares y Observaciones Solares a Largo Plazo}, foreign = {Sunspot Index and Long-term Solar Observations}}
\DeclareAcronym{SMEI}{short = SMEI, long = {Generador de Imágenes de Eyecciones de Masa Solar}, foreign = {Solar Mass Ejection Imager}}
\DeclareAcronym{SOHO}{short = SOHO, long = {Observatorio Solar y Heliosférico}, foreign = {Solar and Heliospheric Observatory}}
\DeclareAcronym{SOLA}{short = SOLA, long = {Promedios Localizados Óptimamente Sustractivos}, foreign = {Subtractive Optimally Localized Averages}}
\DeclareAcronym{STEREO}{short = STEREO, long = {Observatorio de Relaciones Solares-Terrestres}, foreign = {Solar TErrestrial RElations Observatory}}
\DeclareAcronym{Sunspots}{short = Sunspots, long = {Manchas Solares}}
\DeclareAcronym{SW}{short = SW, long = {Viento Solar}, foreign = {Solar Wind}}
\DeclareAcronym{SWEAP}{short = SWEAP, long = {Electrones, Alfas y Protones del Viento Solar}, foreign = {Solar Wind Electrons Alphas and Protons}}
\DeclareAcronym{SWOOPS}{short = SWOOPS, long = {Observaciones del Viento Solar Sobre los Polos del Sol}, foreign = {Solar Wind Observations Over the Poles of the Sun}}
\DeclareAcronym{TRACE}{short = TRACE, long = {Explorador de la Región de Transición y Coronales}, foreign = {Transition Region and Coronal Explorer}}
\DeclareAcronym{UVCS}{short = UVCS, long = {Espectrómetro Coronógrafo Ultravioleta}, foreign = {Ultraviolet Coronagraph Spectrometer}}
\DeclareAcronym{WDC}{short = WDC, long = {Centro Mundial de Datos}, foreign = {World Data Center}}
\DeclareAcronym{WFI}{short = WFI, long = {Generador de Imágenes de Campo Amplio}, foreign = {Wide Field Imager}}
\DeclareAcronym{WISPR}{short = WISPR, long = {Generador de Imágenes de Campo Amplio para Solar Probe}, foreign = {Wide-field Imager for Solar Probe}}
 
